\markboth{摘要}{}
\noindent{\sffamily\fontsize{38}{44}\selectfont\bfseries\color{navy}EAR}\hfill{\sffamily\footnotesize\color{muted}中文技术报告\quad /\quad 2026.10}\par
\vspace{1mm}\noindent{\sffamily\small\color{muted}Effective Auto Research}\par
\vspace{4mm}\noindent{\sffamily\fontsize{20}{28}\selectfont\bfseries\color{ink}以反馈推进研究，\\以研究结果改进策略}\par
\vspace{4mm}\noindent\AuthorsCompact
\vspace{3mm}\noindent\textcolor{navy!35}{\rule{\textwidth}{.5pt}}\par
\vspace{1mm}\subsection*{摘要}
EAR（Effective Auto Research）是一套开源科研自动化系统，面向问题发现、自主数学研究与评审回应。它以减少研究者的主动投入为目标，同时将研究质量要求与机器预算作为约束。科研往往在计划改变时向前推进：新文献改变选题判断，反例迫使推导转向，评阅将稿件送回修订。同类工作已将自动化带入选题、实验与写作；从这些环节走向持续研究，还需要让材料能够接续、反馈能够指导行动，并让一次任务的经验进入下一次执行。

EAR 的核心理念是\textbf{以反馈推进研究，以研究结果改进策略}。AutoVibeIdea 通过文献检索、批判与树搜索，将研究方向收敛为可执行提案；AutonomousMath 的内循环推进推导、验证与稿件修订，外循环用完整研究结果检验并演化执行策略；AutoRebuttal 从审稿疑虑与已有证据出发，补齐缺口、比较并修订回应。三个模块可独立使用，也可由研究者通过提案、稿件与评阅材料接续不同阶段的工作。方向选择、关键技术验收与最终审阅由研究者掌握。\src{1}

历史评估中，持续修订将同一组 28 篇稿件的当前版本模型评阅通过数从 3 增至 12；所选演化策略较初始策略获得更多通过结果，每个通过结果所需的研究调用减少 52.8\%。这些结果为 EAR 持续推进研究、改进执行策略的设计提供了初步支持。\src{10}

\vspace{3mm}
\noindent{\small\sffamily\textbf{关键词}\quad 科研自动化；主动人时；批判性树搜索；研究策略演化；证据驱动回应}\par
\clearpage
\markboth{目录}{}
\noindent{\sffamily\small\color{teal}CONTENTS}\par
\vspace{2mm}
\tableofcontents
\thispagestyle{fancy}
\vspace{8mm}
\begin{scope}
\textbf{阅读导航}\quad 第\ref{sec:introduction} 章介绍 EAR 的研究对象与核心理念；第\ref{sec:methods} 章展开三个工作流，重点解释数学研究的双循环；第\ref{sec:evaluation} 章呈现稿件修订、策略比较和回应评分预测的结果及评价条件。
\end{scope}
\clearpage
